\documentclass[12pt,a4paper]{article}
\usepackage[T1]{fontenc}
\usepackage{amsmath,amssymb,mathrsfs,tikz,times,pifont}
\usepackage{enumitem}
\usepackage{multicol}
\usepackage{lmodern}
\usetikzlibrary{trees}
\newcommand\circitem[1]{%
\tikz[baseline=(char.base)]{
\node[circle,draw=gray, fill=red!55,
minimum size=1.2em,inner sep=0] (char) {#1};}}
\newcommand\boxitem[1]{%
\tikz[baseline=(char.base)]{
\node[fill=cyan,
minimum size=1.2em,inner sep=0] (char) {#1};}}
\setlist[enumerate,1]{label=\protect\circitem{\arabic*}}
\setlist[enumerate,2]{label=\protect\boxitem{\alph*}}
%%%::::::by chnini ameur :::::::%%%
\everymath{\displaystyle}
\usepackage[left=1cm,right=1cm,top=1cm,bottom=1.7cm]{geometry}
\usepackage[colorlinks=true, linkcolor=blue, urlcolor=blue, citecolor=blue]{hyperref}
\usepackage{array,multirow}
\usepackage[most]{tcolorbox}
\usepackage{varwidth}
\usepackage{float} %pour utiliser l'option [H] qui force l'image à apparaître exactement à l'endroit où elle est placée dans le code.
\tcbuselibrary{skins,hooks}
\usetikzlibrary{patterns}
%%%::::::by chnini ameur :::::::%%%
\newtcolorbox{exa}[2][]{enhanced,breakable,before skip=2mm,after skip=5mm,
colback=yellow!20!white,colframe=black!20!blue,boxrule=0.5mm,
attach boxed title to top left ={xshift=0.6cm,yshift*=1mm-\tcboxedtitleheight},
fonttitle=\bfseries,
title={#2},#1,
% varwidth boxed title*=-3cm,
boxed title style={frame code={
\path[fill=tcbcolback!30!black]
([yshift=-1mm,xshift=-1mm]frame.north west)
arc[start angle=0,end angle=180,radius=1mm]
([yshift=-1mm,xshift=1mm]frame.north east)
arc[start angle=180,end angle=0,radius=1mm];
\path[left color=tcbcolback!60!black,right color = tcbcolback!60!black,
middle color = tcbcolback!80!black]
([xshift=-2mm]frame.north west) -- ([xshift=2mm]frame.north east)
[rounded corners=1mm]-- ([xshift=1mm,yshift=-1mm]frame.north east)
-- (frame.south east) -- (frame.south west)
-- ([xshift=-1mm,yshift=-1mm]frame.north west)
[sharp corners]-- cycle;
},interior engine=empty,
},interior style={top color=yellow!5}}
%%%%%%%%%%%%%%%%%%%%%%%
\usepackage{fancyhdr}
\usepackage{eso-pic}         % Pour ajouter des éléments en arrière-plan
% Commande pour ajouter du texte en arrière-plan
\usepackage{tkz-tab}
\AddToShipoutPicture{
    \AtTextCenter{%
        \makebox[0pt]{\rotatebox{80}{\textcolor[gray]{0.7}{\fontsize{5cm}{5cm}\selectfont PGB}}}
    }
}
\usepackage{lastpage}
\fancyhf{}
\pagestyle{fancy}
\renewcommand{\footrulewidth}{1pt}
\renewcommand{\headrulewidth}{0pt}
\renewcommand{\footruleskip}{10pt}
\fancyfoot[R]{
\color{blue}\ding{45}\ \textbf{2025}
}
\fancyfoot[L]{
\color{blue}\ding{45}\ \textbf{Prof:M. BA}
}
\cfoot{\bf
\thepage /
\pageref{LastPage}}
% Création du compteur pour les exercices
\newcounter{exercice}
\renewcommand{\theexercice}{\arabic{exercice}}  % Définit l'affichage du compteur en chiffres arabes

% Définir la commande \exo
\newcommand{\exo}{\refstepcounter{exercice}\textbf{Exercice \theexercice} }
\begin{document}
\renewcommand{\arraystretch}{1.5}
\renewcommand{\arrayrulewidth}{1.2pt}
\begin{tikzpicture}[overlay,remember picture]
    \node[draw=blue,line width=1.2pt,fill=purple,text=blue,inner sep=3mm,rounded corners,pattern=dots]at ([yshift=-2.5cm]current page.north) {\begingroup\setlength{\fboxsep}{0pt}\colorbox{white}{\begin{tabular}{|*1{>{\centering \arraybackslash}p{0.28\textwidth}} |*2{>{\centering \arraybackslash}p{0.2\textwidth}|} *1{>{\centering \arraybackslash}p{0.19\textwidth}|} }
                \hline
                \multicolumn{3}{|c|}{$\diamond$$\diamond$$\diamond$\ \textbf{Lycée de Dindéfélo}\ $\diamond$$\diamond$$\diamond$ } & \textbf{A.S. : 2026/2027}                                              \\ \hline
                \textbf{Matière: Mathématiques}                                                                                    & \textbf{Niveau : }\textbf{TL} & \textbf{Date: 05/11/2026} & \textbf{} \\ \hline
                \multicolumn{4}{|c|}{\parbox[c]{10cm}{\begin{center}
                                                                  \textbf{{\Large\sffamily Composition Des Applications}}
                                                              \end{center}}}                                                                                                        \\ \hline
            \end{tabular}}\endgroup};
\end{tikzpicture}

\vspace{3cm}
\fbox{\textbf{\exo}}

\begin{enumerate}
    \item \textbf{Réponse : a)} $-\sqrt{x^2 + 1}$
    \item \textbf{Réponse : a)} $g \circ f(x) = x^2 + 2x$
    \item \textbf{Réponse : a)} $f \circ g(x) = \frac{1}{x-3}$
    \item \textbf{Réponse : a)} $(f \circ g)(x) = \frac{x^2 - 3x + 3}{x + 2}$
    \item \textbf{Réponse : b)} $g \circ f(x) = x^2 + 2x$
\end{enumerate}

\vspace{0.3cm}

\fbox{\textbf{\exo}}
\begin{itemize}
    \item \textbf{Faux}. La composition des applications n'est pas commutative en général (c'est-à-dire que $f \circ g \neq g \circ f$).
\end{itemize}

\vspace{0.3cm}

\fbox{\textbf{\exo}}
\begin{enumerate}
    \item \textbf{Domaines de définition :} $D_f = \mathbb{R}$ et $D_g = [3, +\infty[$.
    \item \textbf{Expressions :} 
    
    $f \circ g(x) = f(\sqrt{x-3}) = 2(\sqrt{x-3})^2 + 1 = 2(x-3) + 1 = 2x - 5$ (pour $x \ge 3$).
    
    $g \circ f(x) = g(2x^2 + 1) = \sqrt{(2x^2 + 1) - 3} = \sqrt{2x^2 - 2}$ (pour $2x^2 - 2 \ge 0$, soit $x \in ]-\infty, -1] \cup [1, +\infty[$).
    \item \textbf{Calculs :} 
    
    $f \circ g(3) = 2(3) - 5 = 1$.
    
    $g \circ f(-2) = \sqrt{2(-2)^2 - 2} = \sqrt{8 - 2} = \sqrt{6}$.
\end{enumerate}

\vspace{0.3cm}

\fbox{\textbf{\exo}}
\begin{enumerate}
    \item \textbf{Calculs des composées :}
    
    $g \circ f(x) = g(-x^2 + 3x) = \frac{(-x^2 + 3x) + 1}{(-x^2 + 3x) - 2} = \frac{-x^2 + 3x + 1}{-x^2 + 3x - 2}$.
    
    $f \circ g(x) = f\left(\frac{x+1}{x-2}\right) = -\left(\frac{x+1}{x-2}\right)^2 + 3\left(\frac{x+1}{x-2}\right)$.
    \item \textbf{Image de $-1$ par $g \circ f$ (deux manières) :}
    \begin{itemize}
        \item \textit{1ère méthode :} On calcule d'abord $f(-1) = -(-1)^2 + 3(-1) = -1 - 3 = -4$, puis on calcule $g(-4) = \frac{-4 + 1}{-4 - 2} = \frac{-3}{-6} = \frac{1}{2}$.
        \item \textit{2ème méthode :} On utilise l'expression générale de $g \circ f(x)$ en remplaçant $x$ par $-1$ : \\
        $g \circ f(-1) = \frac{-(-1)^2 + 3(-1) + 1}{-(-1)^2 + 3(-1) - 2} = \frac{-1 - 3 + 1}{-1 - 3 - 2} = \frac{-3}{-6} = \frac{1}{2}$.
    \end{itemize}
\end{enumerate}

\vspace{0.3cm}

\fbox{\textbf{\exo}}
\begin{enumerate}
    \item \textbf{Domaines :} $D_f = \mathbb{R}$, $D_g = \mathbb{R}$, et $D_h = [2, +\infty[$.
    \item \textbf{Calcul de $f \circ g(x)$ :}
    $f \circ g(x) = f(2x - 3) = (2x - 3)^2 - 2(2x - 3) + 2 = 4x^2 - 12x + 9 - 4x + 6 + 2 = 4x^2 - 16x + 17$.
    \item \textbf{Calcul de $h \circ f(x)$ :}
    $h \circ f(x) = h(x^2 - 2x + 2) = \sqrt{(x^2 - 2x + 2) - 2} = \sqrt{x^2 - 2x} = \sqrt{x(x-2)}$.
\end{enumerate}

\vspace{0.3cm}

\fbox{\textbf{\exo}}
\begin{enumerate}
    \item[a)] \textbf{Calculs :}
    $f(-1) = \frac{-(-1)}{(-1)^2 + 3} = \frac{1}{4}$. Puis $f \circ f(-1) = f\left(\frac{1}{4}\right) = \frac{-\frac{1}{4}}{\left(\frac{1}{4}\right)^2 + 3} = \frac{-\frac{1}{4}}{\frac{1}{16} + 3} = \frac{-\frac{1}{4}}{\frac{49}{16}} = -\frac{4}{49}$.
    $g \circ f(2) = g\left(\frac{-2}{2^2 + 3}\right) = g\left(-\frac{2}{7}\right) = 1 - \left(-\frac{2}{7}\right)^2 = 1 - \frac{4}{49} = \frac{45}{49}$.
    \item[b)] \textbf{Expressions :}
    
    $f \circ g(x) = f(1 - x^2) = \frac{-(1 - x^2)}{(1 - x^2)^2 + 3} = \frac{x^2 - 1}{x^4 - 2x^2 + 4}$.
    
    $g \circ g(x) = g(1 - x^2) = 1 - (1 - x^2)^2 = 1 - (1 - 2x^2 + x^4) = 2x^2 - x^4$.
\end{enumerate}

\vspace{0.3cm}

\fbox{\textbf{\exo}}
\begin{itemize}
    \item \textbf{Calculs :}
    
    $g \circ f(x) = g(2x - 3) = \frac{-3(2x - 3) + 7}{(2x - 3) + 2} = \frac{-6x + 9 + 7}{2x - 1} = \frac{-6x + 16}{2x - 1}$.
    
    $h \circ f(x) = h(2x - 3) = \sqrt{4(2x - 3) - 1} = \sqrt{8x - 12 - 1} = \sqrt{8x - 13}$.
    
    $f \circ g(x) = f\left(\frac{-3x+7}{x+2}\right) = 2\left(\frac{-3x+7}{x+2}\right) - 3 = \frac{-6x + 14 - 3(x + 2)}{x + 2} = \frac{-9x + 8}{x + 2}$.
\end{itemize}

\vspace{0.3cm}

\fbox{\textbf{\exo}}
\begin{enumerate}
    \item \textbf{Calculs :}
    
    $f \circ g(x) = f(\sqrt{4x - 1}) = 2(\sqrt{4x - 1}) - 3$.
    
    $g \circ f(x) = g(2x - 3) = \sqrt{4(2x - 3) - 1} = \sqrt{8x - 13}$.
    \item \textbf{Détermination de $p$ et $q$ pour $h(x) = \frac{x^2 - 2}{x^2 + 2}$ :}
    
    En posant $p(x) = x^2$ et $q(x) = \frac{x - 2}{x + 2}$, on a bien $(q \circ p)(x) = q(x^2) = \frac{x^2 - 2}{x^2 + 2} = h(x)$.
\end{enumerate}

\vspace{0.3cm}

\fbox{\textbf{\exo}}
\begin{itemize}
    \item \textbf{Réponse :} $(g \circ f)(x) = g(2x + 3) = 6(2x + 3) - 18 = 12x + 18 - 18 = 12x$. \textbf{Option a)}.
\end{itemize}

\vspace{0.3cm}

\fbox{\textbf{\exo}}
\begin{enumerate}
    \item \textbf{Vérification :} 
    $g(2) = 3(2) - 4 = 2$. Puis $f(g(2)) = f(2) = 2(2) + 3 = 7$. Donc $(f \circ g)(2) = 7$.
    \item \textbf{Calcul :} 
    $f\left(\frac{-3}{2}\right) = 2\left(\frac{-3}{2}\right) + 3 = -3 + 3 = 0$.
    Puis $(g \circ f)\left(\frac{-3}{2}\right) = g(0) = 3(0) - 4 = -4$.
\end{enumerate}

\vspace{0.3cm}

\fbox{\textbf{\exo}}
\begin{enumerate}
    \item \textbf{Expressions :}
    $f \circ g(x) = f(2x + 1) = -2(2x + 1)^2 + 7 = -2(4x^2 + 4x + 1) + 7 = -8x^2 - 8x + 5$.
    $g \circ f(x) = g(-2x^2 + 7) = 2(-2x^2 + 7) + 1 = -4x^2 + 14 + 1 = -4x^2 + 15$.
    \item \textbf{Calcul de $g \circ f(2)$ de deux manières :}
    \begin{itemize}
        \item \textit{1ère méthode :} $f(2) = -2(2)^2 + 7 = -8 + 7 = -1$, puis $g(-1) = 2(-1) + 1 = -1$.
        \item \textit{2ème méthode :} En utilisant l'expression de $g \circ f(x)$ : \\
        $g \circ f(2) = -4(2)^2 + 15 = -4(4) + 15 = -16 + 15 = -1$.
    \end{itemize}
\end{enumerate}

\end{document}